%% ============================================================
%%  Front matter — Abstract / Resumen / Resum
%%
%%  Author      : Carlos David Prado-Socorro
%%  Institution : ICMol, Universitat de València
%%
%%  Trilingual abstract: English (main), Spanish (Resumen) and
%%  Valencian/Catalan (Resum). The Spanish and Catalan pages are
%%  wrapped in babel otherlanguage environments so they hyphenate
%%  correctly; English remains the document's main language.
%% ============================================================

%% ---- Standalone preamble (skipped when \include'd from thesis.tex) ----
\ifdefined\thesismode\else
  \documentclass[12pt, twoside]{report}
  \usepackage{thesis-format}
  \addbibresource{../bibliography/references.bib}
  \begin{document}
\fi

%% ============================================================
\chapter*{Abstract}
\addcontentsline{toc}{chapter}{Abstract}
\markboth{Abstract}{Abstract}
%% ============================================================

The energy and latency of modern computing are increasingly dominated by the
movement of data between separate memory and processing units rather than by
arithmetic itself. One physical response is to compute with devices whose
internal dynamics carry and transform information, as the biological synapse
does. This thesis develops a \emph{volatile organic memristive device} (a
two-terminal element made from an ion-containing polymer--electrolyte composite)
and argues that its defining property, a conductance that retains a decaying
record of recent electrical history, is a computational resource rather than a
shortcoming: a controlled \emph{fading memory} that makes the device a physical
leaky integrator for temporal computing.

The thesis first demonstrates the complete synaptic measurement suite in a
proof-of-concept Super~Yellow/Hybrane/lithium-triflate device: current--voltage
hysteresis, potentiation and depression, an excitatory post-synaptic current,
distinguishable short- and long-term retention, and spike-timing-dependent
plasticity. These responses are interpreted through the mobility asymmetry
between the triflate anion and the lithium cation. The fabrication archive then
reveals a campaign-scale reproducibility collapse of the Hybrane line. Changing
to a poly(ethylene oxide) host resolves the collapse and leads to a
fabrication-record methodology and an automated feature pipeline. On the PEO
platform, a replicated composition series tunes the fading-memory time over
roughly 2--20\,s. Host, anion and cation substitutions remain a sample-limited
survey, and the data do not support the hypothesised \ce{Li+} > \ce{Na+} >
\ce{K+} ordering. Finally, compact behavioural models fitted to the measured
dynamics drive reservoir-computing simulations. A heterogeneous composition
bank raises the linear memory capacity by a factor of $1.54$ and the
information-processing capacity from $7.1$ to $11.6$, classifies stress on the
WESAD corpus (leave-one-subject-out macro-F1 of $0.894$ binary and $0.758$
three-class), and improves continuous valence/arousal tracking on the CASE
corpus (a $+0.176$ gain in concordance correlation from the spread of
timescales).

Together, the results show that composition tunes fading memory in the
replicated PEO/lithium-triflate series, while mixed-composition banks broaden
temporal memory in simulation. These devices are therefore better compared with
temporal and event-driven neuromorphic hardware than with non-volatile memories.

%% ============================================================
\begin{otherlanguage}{spanish}
\chapter*{Resumen}
\addcontentsline{toc}{chapter}{Resumen}
\markboth{Resumen}{Resumen}
%% ============================================================

El consumo energético y la latencia de la computación moderna están cada vez más
dominados por el movimiento de datos entre unidades de memoria y de
procesamiento separadas, más que por la propia aritmética. Una respuesta física
consiste en computar con dispositivos cuya dinámica interna transporta y
transforma la información, como hace la sinapsis biológica. Esta tesis desarrolla
un \emph{dispositivo memristivo orgánico volátil} (un elemento de dos terminales
fabricado a partir de un compósito polímero--electrolito con iones móviles) y
sostiene que su propiedad definitoria, una conductancia que conserva un registro
decreciente de su historia eléctrica reciente, no es un defecto sino un recurso
computacional: una \emph{memoria evanescente} controlada que convierte al
dispositivo en un integrador con fugas físico para la computación temporal.

La tesis demuestra primero el conjunto completo de medidas sinápticas en un
dispositivo de prueba de concepto Super~Yellow/Hybrane/triflato de litio:
histéresis corriente--voltaje, potenciación y depresión, corriente postsináptica
excitatoria, retención a corto y largo plazo distinguibles, y plasticidad
dependiente del tiempo de disparo. Estas respuestas se interpretan mediante la
asimetría de movilidad entre el anión triflato y el catión litio. A continuación,
el archivo de fabricación revela un colapso de reproducibilidad a escala de
campaña en la línea de Hybrane. El cambio a una matriz de poli(óxido de etileno)
resuelve el colapso y da lugar a una metodología de registro de fabricación y a
una canalización automatizada de extracción de características. En la plataforma
de PEO, una serie de composiciones replicada ajusta el tiempo de memoria
evanescente entre aproximadamente 2 y 20\,s. Las sustituciones de matriz, anión y
catión siguen siendo un estudio limitado por el tamaño muestral, y los datos no
respaldan el orden catiónico hipotético \ce{Li+} > \ce{Na+} > \ce{K+}. Por
último, modelos de comportamiento compactos ajustados a la dinámica medida
alimentan simulaciones de computación de reservorio. Un banco heterogéneo de
composiciones aumenta la capacidad de memoria lineal en un factor de $1{,}54$ y
la capacidad de procesamiento de información de $7{,}1$ a $11{,}6$, clasifica el
estrés en el corpus WESAD (macro-F1 de $0{,}894$ binario y $0{,}758$ de tres
clases con validación dejando-un-sujeto-fuera) y mejora el seguimiento continuo
de valencia/activación en el corpus CASE (una ganancia de $+0{,}176$ en
correlación de concordancia procedente de la diversidad de escalas temporales).

En conjunto, los resultados muestran que la composición ajusta la memoria
evanescente en la serie replicada PEO/triflato de litio, mientras que los bancos
de composiciones mixtas amplían la memoria temporal en simulación. Por ello,
estos dispositivos se comparan mejor con hardware neuromórfico temporal y
dirigido por eventos que con memorias no volátiles.

\end{otherlanguage}

%% ============================================================
\begin{otherlanguage}{catalan}
\chapter*{Resum}
\addcontentsline{toc}{chapter}{Resum}
\markboth{Resum}{Resum}
%% ============================================================

El consum energètic i la latència de la computació moderna estan cada vegada més
dominats pel moviment de dades entre unitats de memòria i de processament
separades, més que per la mateixa aritmètica. Una resposta física consisteix a
computar amb dispositius la dinàmica interna dels quals transporta i transforma
la informació, com fa la sinapsi biològica. Aquesta tesi desenvolupa un
\emph{dispositiu memristiu orgànic volàtil} (un element de dos terminals fabricat
a partir d'un compòsit polímer--electròlit amb ions mòbils) i defensa que la seua
propietat definitòria, una conductància que conserva un registre decreixent de la
seua història elèctrica recent, no és un defecte sinó un recurs computacional:
una \emph{memòria evanescent} controlada que converteix el dispositiu en un
integrador amb fugues físic per a la computació temporal.

La tesi demostra primer el conjunt complet de mesures sinàptiques en un
dispositiu de prova de concepte Super~Yellow/Hybrane/triflat de liti: histèresi
corrent--voltatge, potenciació i depressió, corrent postsinàptic excitatori,
retenció a curt i llarg termini distingibles, i plasticitat dependent del temps
de dispar. Aquestes respostes s'interpreten mitjançant l'asimetria de mobilitat
entre l'anió triflat i el catió liti. A continuació, l'arxiu de fabricació revela
un col·lapse de reproductibilitat a escala de campanya en la línia de Hybrane. El
canvi a una matriu de poli(òxid d'etilè) resol el col·lapse i dóna lloc a una
metodologia de registre de fabricació i a una canalització automatitzada
d'extracció de característiques. En la plataforma de PEO, una sèrie de
composicions replicada ajusta el temps de memòria evanescent entre aproximadament
2 i 20\,s. Les substitucions de matriu, anió i catió continuen sent un estudi
limitat per la grandària mostral, i les dades no confirmen l'ordre catiònic
hipotètic \ce{Li+} > \ce{Na+} > \ce{K+}. Finalment, models de comportament
compactes ajustats a la dinàmica mesurada alimenten simulacions de computació de
reservori. Un banc heterogeni de composicions augmenta la capacitat de memòria
lineal en un factor de $1{,}54$ i la capacitat de processament d'informació de
$7{,}1$ a $11{,}6$, classifica l'estrès en el corpus WESAD (macro-F1 de $0{,}894$
binari i $0{,}758$ de tres classes amb validació deixant-un-subjecte-fora) i
millora el seguiment continu de valència/activació en el corpus CASE (un guany de
$+0{,}176$ en correlació de concordança procedent de la diversitat d'escales
temporals).

En conjunt, els resultats mostren que la composició ajusta la memòria evanescent
en la sèrie replicada PEO/triflat de liti, mentre que els bancs de composicions
mixtes amplien la memòria temporal en simulació. Per tant, aquests dispositius es
comparen millor amb maquinari neuromòrfic temporal i dirigit per esdeveniments
que amb memòries no volàtils.

\end{otherlanguage}

%% ---- Standalone close (skipped when \include'd from thesis.tex) ----
\ifdefined\thesismode
  \let\thesisstandaloneclose\relax
\else
  \def\thesisstandaloneclose{\end{document}}
\fi
\thesisstandaloneclose
